%==============================================================================
% APÊNDICE B — REFERÊNCIAS ADICIONAIS POR CATEGORIA
%==============================================================================
\chapter{Referências Adicionais por Categoria}
\label{apendice:referencias}

Este apêndice lista referências complementares organizadas por tema,
que podem ser consultadas para aprofundamento.

\section{Processamento de Imagens e Deep Learning}

\begin{itemize}
  \item \cite{lecun2015deep} — A obra seminal sobre deep learning.
  \item \cite{ronneberger2015u} — O artigo original da arquitetura U-Net.
  \item \cite{goodfellow2016deep} — O livro-texto definitivo sobre deep learning.
\end{itemize}

\section{IA na Odontologia}

\begin{itemize}
  \item \cite{arnet2022cancer} — Revisão sistemática sobre DL em câncer oral.
  \item \cite{oliveira2020lesoes} — Classificação de lesões bucais com ResNet-50.
  \item \cite{silva2017caries} — Detecção automatizada de cáries.
\end{itemize}

\section{Gêmeos Digitais e Novas Fronteiras}

\begin{itemize}
  \item \cite{joda2023twins} — Gêmeos digitais em odontologia protética.
  \item \cite{claro2024framework} — Framework SDD/TDD para sistemas odontológicos.
\end{itemize}

\section{Tutoriais e Implementações}

\begin{itemize}
  \item Documentação oficial do TensorFlow: \url{https://www.tensorflow.org/}
  \item Documentação do PyTorch: \url{https://pytorch.org/}
  \item Google Colab: \url{https://colab.research.google.com/}
  \item OpenCV: \url{https://opencv.org/}
  \item Scikit-learn: \url{https://scikit-learn.org/}
\end{itemize}